\section{Introduction}

Agentic assistants for electronic design automation (EDA) are already deployed in-house, coordinating design flows through Flash-tier masters.
% 中文：面向电子设计自动化（EDA）的智能体助手已在公司内部署，通过 Flash 档 master 协调设计流程。
Yet a correctness-critical slice of the flow resists them: in custom and analog design, the signal that determines whether a layout ``works'' lives in \emph{visual} artifacts---layout, schematic, and waveform views---that a text-only model cannot observe.
% 中文：但流程中与正确性紧密相关的一部分却难以被它们覆盖：在定制与模拟设计中，决定版图是否「可用」的信号存在于版图、原理图与波形视图这类视觉产物中，而纯文本模型无法观测。
A subagent that can only read scripts and logs is blind to the very object it is designing, and must rely on an engineer to translate geometry into text---a bottleneck that scales with design volume.
% 中文：一个只能读取脚本与日志的子智能体，对其正在设计的对象本身是盲的，必须依赖工程师把几何信息翻译成文本——这一瓶颈随设计规模增长。

\paragraph{Contributions.}
We build the multimodal core of \textbf{BaiZe}---an EDA agent---on top of the BaiZe language backbone~\cite{baiZe-llm}, and ask a single question: \emph{can a from-scratch vision encoder make such a small model see the visual EDA artifacts that carry the correctness signal?}
% 中文：我们把 \textbf{BaiZe}——一个高性价比 EDA 智能体——的多模态内核构建在 BaiZe 语言骨干 \cite{baiZe-llm} 之上，并回答一个核心问题：\emph{一个从零训练的视觉编码器，能否让这样一个小模型「看见」承载正确性信号的视觉 EDA 产物？}

\begin{itemize}
\item \textbf{A from-scratch vision encoder.}
BaiZe pre-trains its own vision encoder rather than reusing an off-the-shelf one, and selects the tower from iso-parameter candidates at the $500$--$600$\,M scale.
% 中文：从零训练的视觉编码器。BaiZe 自研而非复用现成的视觉编码器，并在 $500$–$600$M 参数量级的同参数候选中选出编码器塔。
A ViT-L/14 tower reaches $68.2\%$ linear-probe on an internal protocol, $+3.9$\,pp above ViT-B/16, at a $196$-token-per-image budget.
% 中文：ViT-L/14 塔在内部协议上线性探测达 $68.2\%$，比 ViT-B/16 高 $+3.9$\,pp，且每图仅 196 token 预算。

\item \textbf{Token-bounded multimodal alignment.}
BaiZe connects the encoder to the frozen language backbone through a lightweight projector, following the staged curriculum of open multimodal assistants, and asks whether a $196$-token-per-image budget suffices for detail-dense layout and schematic views.
% 中文：成本受限的多模态对齐。BaiZe 通过轻量投影层把编码器连接到冻结的语言骨干，沿用开放多模态助手的分阶段课程，并检验每图 196 token 的预算是否足以应对细节密集的版图与原理图视图。
We show that encoder-quality gains do \emph{not} transfer downstream: the larger tower averages $26.1\%$ over 11 benchmarks, indistinguishable from the smaller one, and we localise the gap to the $224$-px single-tile setting and the alignment architecture.
% 中文：我们发现编码器质量的增益并未迁移到下游：更大的塔在 11 个基准上平均 $26.1\%$，与小塔不可分辨；我们把落差定位到 224 像素单块设定与对齐架构。

\item \textbf{Domain adaptation.}
A final reinforcement-learning stage adapts the multimodal model to proprietary semiconductor practice, learning team conventions that no public benchmark measures.
% 中文：领域适配。最后一个强化学习阶段把多模态模型适配到专有半导体实践，学习任何公开基准都无从衡量的团队约定。

\item \textbf{Three-tier evaluation.}
We evaluate end-to-end on 12 general multimodal benchmarks, 4 domain (EDA) benchmarks, and the multimodal subset of an EDA flow benchmark, under a single frozen execution-based protocol.
% 中文：三层评测。我们在 12 个通用多模态基准、4 个领域（EDA）基准，以及 EDA 流程基准的多模态子集上，按统一的冻结的、基于执行的协议进行端到端评测。
\end{itemize}